% SPDX-License-Identifier: Apache-2.0
% covers: ScanVideo, ScanTap, VidMux, txhdl_parts::hdmi::Raster
\datasheet{ScanVideo}{The video peripheral with a scanout from memory
beside it, programmed over one AXI-Lite slot.}

\dsfacts{\code{lib/parts/src/scanout.rs}}%
{\code{txhdl\_parts::scanout::ScanVideo<HV, HFP, HSW, HBP, VV, VFP,
VSW, VBP, SHIFT, AW, TOTAL, STRIDE, LO, HI>}}%
{\code{//docs:examples}}

\subsection*{Function}

\code{ScanVideo} is what the flagship's video slot holds (issue 151): the
video peripheral \code{Hdmi}, unchanged, and beside it a scanout that
shows a frame from memory, with a multiplexer that puts one picture or
the other on the pins. It is a unit of seven, on the pixel clock:

\begin{itemize}
\item \code{LiteSplit} divides the slot by address bit 7: \code{Hdmi}'s
words below \code{0x80}, \code{ScanCtl}'s from there.
\item \code{Raster}, in \code{lib/parts/src/hdmi.rs}, counts the beam
again with \code{Hdmi}'s own functions and parameters, for
\code{LinePair}. It drives the column, the visible part, the line and
frame starts and the row from registers. It is the scanout's raster,
not Razboj's rasteriser, which has a sheet of its own under the same
name (issue 1055).
\item \code{LinePair} holds two lines and asks for each one or two
rows ahead, on \code{req}; their words arrive on \code{words}. Both
channels leave the unit for the bus clock. \code{ScanCtl}'s show bit,
which picks the scanout's picture at \code{VidMux}, is also
\code{LinePair}'s \code{show}, so the fetch starts only once the
scanout is shown (issue 1178). \code{ScanCtl}'s ahead bit is
\code{LinePair}'s \code{two} (issue 1523).
\item \code{ScanTap} puts \code{LinePair}'s underflow and stuck bits,
the stuck line's address, the longest line and the count of late
lines on a channel back to \code{ScanCtl}.
\item \code{VidMux} puts \code{Hdmi}'s pixel or the scanout's, a
\code{0x00RRGGBB} word, on the pins, with \code{Hdmi}'s syncs either
way.
\end{itemize}

The children run in an order that puts every wire's driver before its
reader: \code{Raster} and \code{ScanCtl}, then \code{LinePair}, then
\code{ScanTap} and \code{Hdmi}, then \code{VidMux}. The one path the
other way, from \code{LinePair} back to \code{ScanCtl}, is a channel.
Counting the beam twice is what keeps that order possible, and
\code{ex\_scanvideo} checks the two counts agree on every cycle.

\subsection*{Parameters}

\begin{center}\footnotesize
\begin{tabular}{@{}lp{0.68\textwidth}@{}}
\toprule
Parameter & Meaning \\
\midrule
\code{HV} to \code{VBP}, \code{SHIFT} & \code{Hdmi}'s: the video mode, and
a framebuffer pixel's size on the screen. \\
\code{AW} & The width of a column; \code{HV} is at most \code{1 << AW}. \\
\code{TOTAL} & The rows of a frame, \code{VV + VFP + VSW + VBP}; the
\code{Default} refuses one that disagrees. \\
\code{STRIDE} & The bytes from one line to the next in memory. \\
\code{LO}, \code{HI} & The frame's memory, the bytes from \code{LO} up
to but not including \code{HI}, outside which no line is asked for. \\
\bottomrule
\end{tabular}
\end{center}

The tables below use the flagship's: 640 by 480 at 60~Hz, a
framebuffer pixel of four by four, ten bits of column, 525 rows and
4096 bytes a line, as \code{hdmi/src/netlist.rs} writes \code{scan\_video}.
A line in memory is Razboj's row of 1024 words, of which the first 640
are shown (issue 985). The frame's memory is the board's gigabyte from
\code{0x4000\_0000} to \code{0x8000\_0000} (issue 1382).

\dstables{ScanVideo}

\subsection*{Behaviour}

\begin{itemize}
\item The pair's line is \code{HV} words, and its rows are \code{VV}
visible of \code{TOTAL}.
\item A line not wholly inside \code{LO} to \code{HI} is not asked
for: the pair sets its stuck bit with that line's address, the row
shows \code{LATE}, and the frame goes on (issue 1382). A base gone
wrong once sent a line's burst into the serial port's page, and the
core's prints stopped behind it.
\item With the ahead bit set, the pair asks for lines two rows ahead
from the next frame, and a line's words wait in the crossing until
the row before its own (issue 1523).
\end{itemize}

\subsection*{Verification}

\begin{itemize}
\item \code{ex\_scanvideo} runs the whole unit with \code{LineFetch},
\code{ScanFetch} and a memory that stalls on command. It paints a
pixel through \code{Hdmi} and sees it, shows two frames of scanout and
checks every visible pixel, reads the underflow bit over the slot as a
stall sets it and a write clears it, and asserts that \code{Raster}'s
count is \code{Hdmi}'s on every cycle, across a reset.
\item \code{//docs:scanvideo\_sim\_scanvideo\_tb\_test} and
\code{//docs:scanvideo\_vsim\_test} simulate its netlist against that
run under nvc and Verilator.
\item \code{//lib/parts:parts\_test}'s
\code{a\_frame\_with\_its\_words}\allowbreak
\code{\_on\_time\_has\_no\_late\_line} runs the flagship's pair, with this
sheet's \code{LO} and \code{HI}, and raster over three frames, one row
ahead and two, and sees no late line once the first frame is over.
\item \code{//cpu/vreteno:board\_test}'s
\code{a\_base\_outside\_the\_memory}\allowbreak
\code{\_is\_never\_fetched} gives the flagship's line the serial
port's page as its base: no line is asked for, the stuck bit is set at
that base, and the core's greeting comes out whole.
\item On the board, \code{scanprobe} shows a frame from the DDR3 and
reads the underflow bit idle and under load; both read clear
(\code{docs/board-checks.md}, PR 1041's session).
\end{itemize}

\subsection*{Used by}

\code{flagship/board/flagship.v}, in the video slot at \code{0x3200},
with \code{chan\_cdc} transferring \code{req} and \code{words} across
clock domains to the board's \code{ScanFetch} and \code{LineFetch}
engines. The words' crossing, \code{swords\_cdc}, holds 2048 words,
more than two lines, so the bus's read data never waits on the pixel
side, and two rows ahead the next line waits there whole.

\subsection*{Limits}

The module operates in a single video mode determined by its synthesis
parameters, and its frame's memory is fixed the same way. The
scanout's picture is the native video mode's, without scaling. The
bits \code{ScanCtl} reads reach it up to three cycles after
\code{LinePair} drives them, through the channel between the two.
